\begin{lemma}[The terminal stored set lies in the last indexed stored set]
For every hierarchy $K$ with $d$ dividers, its terminal stored set
$D^K_{n,d+1}$ is contained in its last indexed stored set $D^K_{n,d}$.
Moreover, every color in $D^K_{n,d}\setminus D^K_{n,d+1}$ is missing on
the terminal tree.
\end{lemma}
\begin{proof}
Let $R$ be the terminal-rung \noderef{TashkinovStoredColors} set and let
$A$ be the last \noderef{HierarchyStarredStage}.  The hierarchy fields
give $A\preccurlyeq T$, where $T$ is the terminal tree.  Indeed, when
$d=0$, $A$ is the base core and the core-inside-first field of
\noderef{TashkinovHierarchy}, followed by its terminal-stage equality,
gives the claim.  When $d>0$, $A$ is the penultimate copied stage, and the
last nesting field gives the claim directly.

By the entrywise clause of \noderef{TreeSequenceInitialSegment}, every
displayed vertex of $A$ is a displayed vertex of $T$.  Unfolding
\noderef{TreeSequenceVertexSet} therefore gives $V(A)\subseteq V(T)$,
and the union definition in \noderef{TreeMissingColors} gives
$\overline\phi(A)\subseteq\overline\phi(T)$.  Finally unfold
\noderef{HierarchyStoredColorsAt} and
\noderef{HierarchyTerminalStoredColors}:
\[
 D^K_{n,d}=R\setminus\overline\phi(A),\qquad
 D^K_{n,d+1}=R\setminus\overline\phi(T).
\]
If a color belongs to the second difference, it belongs to $R$ and is not
missing on $T$.  The missing-set inclusion implies that it is not missing
on $A$, so it belongs to the first difference.  This proves the stated
containment.  Conversely, if a color belongs to the first difference but
not the second, its membership in the common set $R$ rules out failure of
the second difference through nonmembership in $R$; hence it must be
missing on $T$.  This proves the additional difference clause.
\end{proof}
